\section{Zadanie 11}
\label{set07.task11}
Udowodnij, że jeżeli macierz przejścia $P$ jest symetryczna ($p_{x,y}=p_{y,x}$), to rozkład jednostajny jest rozkładem stacjonarnym dla $P$.

\subsection*{Rozwiązanie}
Skorzystamy z ”ogólnej wskazówki” z zestawu:\\
Niech $P$ będzie macierzą przejścia łańcucha Markowa na zbiorze stanów $S$.
Niech $\pi$ będzie rozkładem prawdopodobieństwa na $S$.
Łatwo uzasadnić, że jeśli dla każdego $x,y\in S$ mamy
\[
\pi_x p_{x,y} = \pi_y p_{y,x}
\]
to $\pi$ jest rozkładem stacjonarnym.

Zauważmy, że $\pi$ – rozkład stacjonarny ($\forall x,y\in S\ \pi_x=\pi_y$) w połączeniu z symetrią
macierzy $P$ implikuje równość, o której mowa we wskazówce: dla każdego $x,y\in S$, zatem rozkład jednostajny jest rozkładem stacjonarnym dla $P$. CKD

